% ============================================================
% SHARED PREAMBLE (COMMON) — texdocs
% Bagian formatting yang TIDAK bergantung ukuran kertas.
% Di-input oleh main.tex SETELAH ../shared/size-a4 atau size-a5.
% Size file menyediakan: geometry, \setstretch, dan macro:
%   \FontChapterHead  \FontSectionHead  \FontKutipan
% ============================================================

% Layout & Font
\usepackage[T1]{fontenc}
\usepackage[utf8]{inputenc}
\usepackage[indonesian]{babel}
\usepackage{mathptmx}          % Times New Roman
\usepackage{indentfirst}

% Heading style
\usepackage{titlesec}

% Gambar
\usepackage{graphicx}
\usepackage{float}
\usepackage{caption}

% Tabel
\usepackage{booktabs}
\usepackage{tabularx}
\usepackage{longtable}
\usepackage{multirow}
\usepackage{array}
\newcolumntype{L}[1]{>{\raggedright\arraybackslash}p{#1}}
\usepackage[table]{xcolor}
\definecolor{headgray}{gray}{0.85}
\usepackage{ragged2e}

% Matematika
\usepackage{amsmath}
\usepackage{amssymb}

% Header/footer
\usepackage{fancyhdr}

% List
\usepackage{enumitem}
\setlist[enumerate,1]{label=\arabic*., leftmargin=*, labelindent=1cm, itemsep=0.3ex, topsep=0.4ex, parsep=0pt}
\setlist[itemize,1]{leftmargin=*, labelindent=1cm, itemsep=0.3ex, topsep=0.4ex, parsep=0pt}

% Daftar isi
\usepackage[titles]{tocloft}

% Hyperlink
\usepackage[hidelinks]{hyperref}

% Daftar pustaka — IEEE style
\usepackage[numbers,sort&compress]{natbib}
\bibliographystyle{IEEEtran}

% Kode program
\usepackage{listings}
\usepackage{color}

% ============================================================
% KONFIGURASI TIPOGRAFI
% ============================================================

% Line spacing (\setstretch) diset di size file.
\setlength{\parindent}{0.5in}
\setlength{\parskip}{0pt}
\setlength{\mathindent}{10mm}

\hyphenpenalty=10000
\exhyphenpenalty=10000
\setlength{\emergencystretch}{3em}

% ============================================================
% FORMAT HEADING (TITLESEC)
% ============================================================

\renewcommand{\chaptername}{BAB}

\titleformat{\chapter}[display]
  {\centering\bfseries\FontChapterHead}
  {\MakeUppercase{\chaptername}\ \thechapter}
  {0pt}
  {\MakeUppercase}
\titlespacing*{\chapter}{0pt}{-20pt}{12pt}

\titleformat{name=\chapter,numberless}[block]
  {\centering\bfseries\FontChapterHead}
  {}
  {0pt}
  {}
\titlespacing*{name=\chapter,numberless}{0pt}{-20pt}{12pt}

\titleformat{\section}
  {\bfseries\FontSectionHead}
  {\makebox[0.5in][l]{\thesection}}
  {0pt}
  {\MakeUppercase}
\titlespacing*{\section}{0pt}{12pt}{6pt}

\titleformat{\subsection}
  {\bfseries\FontSectionHead}
  {\makebox[0.5in][l]{\thesubsection}}
  {0pt}
  {}
\titlespacing*{\subsection}{0pt}{12pt}{6pt}

\setcounter{secnumdepth}{3}
\renewcommand{\thesubsubsection}{\alph{subsubsection}.}
\titleformat{\subsubsection}
  {\normalfont\normalsize}
  {\thesubsubsection}
  {0.5em}
  {}
\titlespacing*{\subsubsection}{0pt}{12pt}{6pt}

% ============================================================
% FORMAT HALAMAN (FANCYHDR)
% ============================================================

\fancypagestyle{mainmatter}{%
  \fancyhf{}%
  \fancyfoot[C]{\thepage}%
  \renewcommand{\headrulewidth}{0pt}%
  \renewcommand{\footrulewidth}{0pt}%
}

\fancypagestyle{plain}{%
  \fancyhf{}%
  \fancyfoot[C]{\thepage}%
  \renewcommand{\headrulewidth}{0pt}%
  \renewcommand{\footrulewidth}{0pt}%
}

% ============================================================
% FORMAT CAPTION
% ============================================================

\captionsetup[figure]{labelsep=period, font=bf, justification=centering}
\captionsetup[table]{labelsep=period, font=bf, justification=centering}

% ============================================================
% PENOMORAN PER BAB
% ============================================================

\numberwithin{equation}{chapter}
\numberwithin{figure}{chapter}
\numberwithin{table}{chapter}

\makeatletter
\def\tagform@#1{\maketag@@@{\bfseries(Persamaan #1)}}
\makeatother

% ============================================================
% NAMA-NAMA INDONESIA
% ============================================================

\renewcommand{\contentsname}{DAFTAR ISI}
\renewcommand{\listfigurename}{DAFTAR GAMBAR}
\renewcommand{\listtablename}{DAFTAR TABEL}
\renewcommand{\bibname}{DAFTAR PUSTAKA}
\addto\captionsindonesian{\renewcommand{\bibname}{DAFTAR PUSTAKA}}
\renewcommand{\figurename}{Gambar}
\renewcommand{\tablename}{Tabel}

% ============================================================
% FORMAT DAFTAR ISI (TOCLOFT)
% ============================================================

\renewcommand{\cftchapfont}{\bfseries}
\renewcommand{\cftchappagefont}{\bfseries}
\renewcommand{\cftchapleader}{\cftdotfill{\cftdotsep}}

\setlength{\cftbeforetoctitleskip}{0pt}
\setlength{\cftaftertoctitleskip}{8pt}
\setlength{\cftbeforeloftitleskip}{0pt}
\setlength{\cftafterloftitleskip}{8pt}
\setlength{\cftbeforelottitleskip}{0pt}
\setlength{\cftafterlottitleskip}{8pt}

% ============================================================
% ENVIRONMENT KUTIPAN PANJANG
% ============================================================

\newenvironment{kutipan}{%
  \begin{list}{}{%
    \setlength{\leftmargin}{10mm}%
    \setlength{\rightmargin}{10mm}%
  }%
  \item\relax\FontKutipan
}{%
  \end{list}
}

% ============================================================
% KONFIGURASI LISTINGS
% ============================================================

\lstset{%
  basicstyle=\ttfamily\small,
  keywordstyle=\bfseries\color{blue!70!black},
  commentstyle=\itshape\color{green!50!black},
  stringstyle=\color{red!60!black},
  numbers=left,
  numberstyle=\tiny\color{gray},
  frame=single,
  breaklines=true,
  captionpos=b
}
